\newif\ifworldscientificclass
\IfFileExists{ws-procs11x85.cls}{\worldscientificclasstrue\documentclass{ws-procs11x85}}{\worldscientificclassfalse\documentclass[11pt]{article}\usepackage[margin=1in]{geometry}}
\IfFileExists{ws-procs-thm.sty}{\usepackage{ws-procs-thm}}{}
\usepackage{graphicx}
\usepackage{placeins}
\usepackage{array}
\setlength{\emergencystretch}{1em}
\graphicspath{{figures/}{paper/figures/}{figures/paper (Long)/}{paper/figures/paper (Long)/}}

% Compatibility fallbacks for local builds with older World Scientific class files.
\providecommand{\keywords}[1]{\par\noindent{\bf Keywords:} #1\par}
\providecommand{\copyrightinfo}[1]{\par\smallskip\noindent{\footnotesize #1}\par\smallskip}
\providecommand{\address}[1]{\par\smallskip\noindent #1\par}
\providecommand{\tbl}[2]{\caption{#1}#2}
\providecommand{\toprule}{\hline}
\providecommand{\colrule}{\hline}
\providecommand{\botrule}{\hline}
\makeatletter
\@ifundefined{itemlist}{\newenvironment{itemlist}{\begin{itemize}}{\end{itemize}}}{}
\makeatother

\begin{document}

\title{GraphGenome-FM: Self-Supervised Topology Pretraining\\
Across Human Pangenome Graphs}

\ifworldscientificclass
\author{Long Nguyen$^{\dag}$ and [Coauthor Names]}

\address{Department of Computer and Information Sciences, Temple University,\\
Philadelphia, PA, USA\\
$^{\dag}$E-mail: [long.nguyen@temple.edu]}

\author{[Principal Investigator Name]}

\address{[Laboratory or Department], Temple University,\\
Philadelphia, PA, USA\\
E-mail: [pi.email@temple.edu]}
\else
\author{Long Nguyen$^{\dag}$ and [Coauthor Names]\\
\small Department of Computer and Information Sciences, Temple University, Philadelphia, PA, USA\\
\small $^{\dag}$E-mail: [long.nguyen@temple.edu]
\and
[Principal Investigator Name]\\
\small [Laboratory or Department], Temple University, Philadelphia, PA, USA\\
\small E-mail: [pi.email@temple.edu]}
\date{}
\maketitle
\fi

\begin{abstract}
Human pangenome graphs encode alternative genomic adjacencies across phased
assemblies, but it remains unclear whether a compact graph encoder can learn
topological representations that transfer across chromosomes and graph
releases. We introduce GraphGenome-FM, a coordinate--topology encoder trained by
query-edge-masked, distance-matched link prediction. We evaluate four versioned
whole-genome graphs: HPRC Release 2 (231 donors, 462 haplotypes), HGSVC3 (65,
130), HPRC Release 1.1, and the official HGSVC3--HPRC-v1 integrated graph. Five
rotating chromosome folds cover chr1--chr22, chrX, and chrY, with three seeds
and two separately defined context tasks. Standalone HPRC models achieve mean
AUPRC 0.950 in strict induced subgraphs and 0.988 in one-hop-expanded
subgraphs; standalone HGSVC3 models achieve 0.975 and 0.987, respectively.
Zero-shot HPRC-to-HGSVC3 transfer is asymmetric (AUPRC 0.600 strict and 0.941
one-hop), and four donors are shared across graph cohorts, so this is
cross-graph rather than strictly independent-donor transfer. Expanded contexts
contain 1.57--1.59 times as many visible nodes and 1.84--1.89 times as many
links, and no interval pairs meet the preregistered 10\% exposure caliper;
strict and expanded results are therefore reported as different tasks, not as a
causal context-improvement comparison. Validation-only temperature scaling
reduces calibration error but does not eliminate it. The current encoder does
not consume nucleotide sequence, aggregate graphs do not support donor-held-out
pretraining, and a chromosome-8 methylation pilot shows negligible incremental
rank gain from graph embeddings beyond strong sequence and annotation features.
These results support transferable pangenome-topology pretraining, but not yet a
general-purpose foundation-model or genome-wide biological-validation claim.
\end{abstract}

\keywords{Pangenome graph; graph neural network; self-supervised learning;
link prediction; HPRC; HGSVC; cCRE; haplotype-specific methylation;
representation transfer.}

% required, do-not-remove
\copyrightinfo{\copyright\ 2026 The Authors. Open Access chapter published by
World Scientific Publishing Company and distributed under the terms of the
Creative Commons Attribution Non-Commercial (CC BY-NC) 4.0 License.}

\section{Introduction}

The human reference genome is moving from a single linear coordinate system
toward graph-based references that represent common and rare variation across
many haplotypes. The Human Pangenome Reference Consortium (HPRC) pangenome and
the complementary Human Genome Structural Variation Consortium (HGSVC)
resources provide such graphs, in which nodes are DNA segments and edges encode
the adjacencies permitted along reference and alternate haplotype
paths.\cite{liao2023hprc,ebert2021hgsvc,garrison2018vg} These graphs make
explicit the structural variation, allelic alternatives, and local topological
context that a single linear sequence cannot express. The human Y chromosome is
an extreme illustration: haplotype-resolved assemblies reveal large recurrent
inversions, palindromic structure, and highly similar segmental duplications whose
arrangement varies extensively between individuals.\cite{hallast2023ychrom} A
linear reference flattens precisely this branching and reconvergence, which a
graph reference is designed to retain.

Most genomic representation learning is nonetheless sequence-first. Linear
models learn strong nucleotide-level representations, but they operate on a
single path through the genome or on a serialized traversal of a
graph.\cite{ji2021dnabert,zhou2024dnabert2,dallatorre2025nt} Serialization
discards exactly the information that makes a pangenome graph useful: branching,
reconvergence, path multiplicity, and the relationship between the reference
backbone and non-reference haplotypes. The graph-native model DeepGene
serializes a Minigraph directed acyclic graph into fixed-length token subgraphs
and applies rotary-position Transformers, but it still abstracts the graph into
a sequence and does not expose edge structure to the
model.\cite{zhang2025deepgene} A representation that aims to be \emph{graph
native} should instead learn from both coordinate continuity and graph topology
at once.

We study self-supervised link prediction as a pretraining objective for human
pangenome graphs. Given a local graph slice and a candidate node pair, the model
predicts whether that edge is present in the pangenome graph. The task is
attractive for three reasons: the labels are intrinsic to the graph and require
no external annotation; the objective scales to unlabeled pangenome data; and
the resulting node embeddings can be reused for downstream functional and
variant annotation. The central question is whether such a graph-native
objective yields representations that generalize across genomic contexts and
across independently constructed pangenome graphs. This paper makes four evidence-bounded
contributions.

\begin{itemlist}
\item \textbf{An all-chromosome HPRC link-prediction benchmark and a single
shared model.} We assemble a benchmark spanning canonical GRCh38 chromosomes with both
strict and one-hop graph closures, and train one shared dual-stream
graph-attention model in place of the many sparse per-window models that an
earlier design required.
\item \textbf{All-chromosome generalization with explicit leakage audits.}
Five rotating folds use every canonical chromosome once as held-out test data,
with query-edge masking, fixed candidate splits, three model seeds, exact slice
exclusions, and validation-only calibration.
\item \textbf{Separate-cohort, pooled-cohort, and release stress tests.} We
evaluate HPRC and HGSVC3 separately, train a shared model on both graphs, test
both transfer directions, and quantify HPRC Release 1.1--2 transfer. Path
metadata identifies the four donors shared by HPRC R2 and HGSVC3, preventing an
independent-donor interpretation.
\item \textbf{An explicit context-exposure audit.} Strict induced-subgraph and
one-hop-expanded prediction are coordinate matched, but no pair is exposure
matched at a 10\% node/link caliper. We therefore treat them as distinct
pretraining tasks and do not claim that larger context is intrinsically
superior.
\item \textbf{A calibrated, reproducible topology benchmark.} We report ranking
metrics as primary outcomes and fit probability temperature and decision
thresholds on validation chromosomes only. Final all-data runs are released as
checkpoint-production artifacts, not generalization estimates.
\end{itemlist}

The strongest current evidence is for transferable local pangenome-topology
learning across chromosomes and graph constructions. Biological transfer,
donor-held-out pretraining, sequence-aware modeling, and validated structural-
variant interpretation remain open requirements rather than current
contributions.

\section{Related Work}

\subsection{Pangenome graphs and haplotype-resolved variation}

Variation graphs represent genomes as paths through shared sequence nodes with
alternate allelic branches. The variation-graph toolkit showed that this
representation improves read mapping by encoding genetic diversity in the
reference itself.\cite{garrison2018vg} Minigraph and the Minigraph--Cactus
pipeline construct reference pangenome graphs in rGFA/GFA form that preserve
linear-reference coordinates while encoding structural
variation.\cite{li2020minigraph,hickey2024minigraphcactus} The HPRC pangenome
applies these ideas at population scale and provides a richer human reference
than GRCh38 alone,\cite{liao2023hprc} while HGSVC contributes
haplotype-resolved assemblies and structural-variation catalogs, including the
complex inversions, palindromes, and segmental duplications of the Y
chromosome, that are valuable for external graph-transfer
evaluation.\cite{ebert2021hgsvc,hallast2023ychrom}

\subsection{Graph neural networks}

Graph neural networks aggregate information over edges and are a general tool
for semi-supervised and self-supervised learning on relational
data.\cite{kipf2017gcn,velickovic2018gat} Graph Attention Networks learn
edge-dependent attention weights, a natural fit for pangenome graphs in which
local topology, branching, and path multiplicity change substantially across
loci.\cite{velickovic2018gat} Our model builds on this family but adds an
explicit coordinate-aware stream, because pangenome nodes also carry genomic
offsets and orientations that a purely topological message-passing scheme
ignores.

\subsection{Genomic foundation models, serialized-graph models, and imputation}

Self-supervised models such as DNABERT, DNABERT-2, and the Nucleotide
Transformer capture useful biological signal from DNA
sequence.\cite{ji2021dnabert,zhou2024dnabert2,dallatorre2025nt} Pangenome graphs
pose a distinct representation problem: the input is not only a string but a
structured graph with multiple valid haplotype paths. DeepGene is the closest
prior work in spirit, using a pangenome graph but serializing it into token
sequences for a Transformer.\cite{zhang2025deepgene} For downstream regulatory
annotation we therefore distinguish three modeling assumptions: linear-reference
models, linearized- or serialized-graph models, and graph-native models that
consume edge structure directly. GraphGenome-FM
belongs to the third class, and we treat the first two as the comparison
baselines that a fair downstream evaluation must include. Our graph-enhancement
application is a form of \emph{imputation}, but it is distinct from genotype
imputation, which infers missing allele values within a reference panel using
linkage disequilibrium and increasingly deep models.\cite{mowlaei2025stici}
GraphGenome-FM instead scores missing \emph{adjacencies} in the graph itself,
which is the natural imputation target once the reference is graph-valued.

\section{Methods}

\subsection{HPRC and HGSVC graph processing}

We derive graph tables from the rGFA/GFA representation of each pangenome. For
every node we record its sequence length, the contig it belongs to, its genomic
offset along that contig, and a reference-status flag indicating whether the node
lies on the linear GRCh38 backbone; edges are directed and carry endpoint
orientation. A positive edge is therefore an observed adjacency in the parsed
pangenome graph, supported by paths in the graph construction, rather than an
experimentally validated functional interaction. HPRC Release 2
(Minigraph--Cactus, GRCh38-aligned) is the primary training and
internal-evaluation graph. HGSVC3 (CHM13-aligned) provides a second-cohort graph;
HPRC Release 1.1 provides a release-transfer control; and the official
HGSVC3--HPRC-v1 graph provides an integrated-graph stress test. All 24 canonical
chromosomes are included, with chrY deliberately retained because of its
structural complexity.\cite{hallast2023ychrom} ENCODE cCRE annotations are
mapped onto HPRC GRCh38 nodes for a downstream pipeline, but no newly executed
cCRE result is part of the frozen evidence set in this manuscript version.
Table~\ref{tab:dataset_summary} summarizes the datasets.

\begin{table}[t]
\tbl{Versioned pangenome graphs in the frozen all-chromosome evidence set.}
{\small\resizebox{\textwidth}{!}{%
\begin{tabular}{@{}lllll@{}}
\toprule
Dataset & Graph build & Reference & Scope & Role\\
\colrule
HPRC & R2 / Minigraph--Cactus SV & GRCh38 & 231 donors / 462 haplotypes & primary\\
HGSVC3 & 2024-02-23 SV graph & CHM13 & 65 donors / 130 haplotypes & cohort transfer\\
HPRC & R1.1 SV graph & GRCh38 & 44 donors / 88 haplotypes & release transfer\\
HGSVC3--HPRC-v1 & official integrated SV graph & CHM13 & 216 haplotypes & stress test\\
\botrule
\end{tabular}}}
\label{tab:dataset_summary}
\end{table}

The exact released GFA, GBZ, README, and checksum-manifest files are retained in
the data ledger. HPRC R1.1 is never pooled with R2 because the releases overlap.
Path metadata verifies four donors shared between HPRC R2 and HGSVC3
(HG00733, HG02818, NA19036, and NA19240); hence cohort transfer is not described
as independent-donor transfer.

\subsection{Graph slice construction and leakage control}

We tile each reference coordinate system deterministically in 5-Mb intervals,
yielding 608 HPRC and 635 HGSVC slices per requested context before loader-level
eligibility filtering. The \emph{strict induced-subgraph} task exposes nodes and
links induced by the reference interval. The \emph{one-hop-expanded} task adds
their immediate graph neighbors. The same coordinate intervals are requested
for both tasks, but exposure is not matched: expanded slices contain 1.57--1.59
times as many visible nodes and 1.84--1.89 times as many visible links, and zero
pairs pass a 10\% node/link exposure caliper. We therefore do not compare the
two task scores as an estimate of the causal benefit of context.

Five rotating folds partition chr1--chr22, chrX, and chrY into non-overlapping
test groups; the next group is used for validation and every remaining
chromosome for training. Each chromosome is a test chromosome exactly once.
The candidate split seed is fixed at 20260806, while three independent model
seeds (42, 314159, and 20260806) quantify initialization variability. Final
all-data checkpoints are trained separately and excluded from generalization
claims.

\subsection{Negative-edge sampling}

Each link-prediction slice contains the positive graph edges together with an
equal number of negative candidate edges. Random non-edges can be trivially
separable when their genomic distance or graph context differs markedly from
that of real edges. We therefore draw distance-matched negatives
(Table~\ref{tab:neg_sampling}): for a positive edge whose endpoints are separated
by an offset $d$, we sample non-edge endpoint pairs from the same sequence
context whose separation lies within $\max(1\,\mathrm{kb},\,0.1\,d)$ of $d$,
optionally also matching endpoint degree. This sampler underpins the robustness
audit reported below.

\begin{table}[t]
\tbl{Negative-sampling controls for link prediction.}
{\small
\begin{tabular}{@{}>{\raggedright\arraybackslash}p{0.25\textwidth}>{\raggedright\arraybackslash}p{0.31\textwidth}>{\raggedright\arraybackslash}p{0.36\textwidth}@{}}
\toprule
Strategy & Non-edge construction & Controlled shortcut and role\\
\colrule
Random &
Sample any absent edge $(u',v')$ uniformly at random from the same slice. &
No matching constraints; useful only as an easy baseline because pairs are often
separable by coordinate distance or graph context.\\
Distance-matched &
For each positive edge $(u,v)$ with endpoint separation $d$, sample absent
pairs whose separation is within $\max(1\,\mathrm{kb},\,0.1\,d)$ of $d$. &
Matches genomic separation and removes the coordinate-distance shortcut; this is
the main hard-negative control, although endpoint degree can still differ.\\
Distance + degree matched &
Within the distance-matched candidate set, prefer absent pairs whose endpoint
degrees match those of the positive edge, using the closest available degrees
when exact matches are unavailable. &
Controls both coordinate distance and endpoint-degree shortcuts; used as an
additional robustness check.\\
\botrule
\end{tabular}}
\label{tab:neg_sampling}
\end{table}

\subsection{Query-edge masking protocol}

For query-edge-masked evaluation, each positive candidate edge is removed from
the message-passing edge index before node embeddings are computed. The model
must therefore score the candidate pair from endpoint features and surrounding
graph context rather than from direct edge availability. Negative candidate
pairs are scored under the same candidate-specific graph construction. The
unmasked setting is retained only as a diagnostic for quantifying how direct
edge availability can inflate link-prediction scores.

\subsection{From per-window to shared training}
\label{sec:shared}

An early version of this work trained one independent model per slice. Those
per-slice models were unstable because each local graph provided too little
positive-edge supervision for reliable strict-closure learning. Pooling slices
into a single self-supervised problem resolves this data-starvation issue: one
shared encoder learns recurring pangenome graph motifs that transfer across
chromosomes. All results below use this shared-training paradigm with
chromosome-level holdout.

\subsection{GraphGenome-FM architecture}

GraphGenome-FM uses a shared dual-stream backbone. Each slice is converted to oriented nodes so
that segment orientation can influence attention and prediction. Raw node
features include sequence length, genomic offset, reference status, graph
degree, local branching fraction, and orientation-derived features. The current
model does not tokenize nucleotide sequence or use k-mer/DNA-language
embeddings; it is a compact coordinate- and topology-pretrained graph
representation model. This distinction matters for cCRE interpretation, because
regulatory-element prediction ultimately requires sequence motifs and assay
signal in addition to graph topology. A node encoder maps these features into a
hidden representation, and two streams then run in parallel:

\begin{itemlist}
\item a \emph{coordinate stream} applies self-attention with multi-scale rotary
position embeddings over genomic offsets, so that dot-product attention depends
on relative genomic distance across local, structural, and chromosomal scales;
\item a \emph{graph stream} applies GAT-style message passing over the observed
edges to capture local topology, branching, and reconvergence; and
\item a \emph{learned fusion gate} combines the two streams per node via a
sigmoid gate $g\,h_{\mathrm{coord}}+(1-g)\,h_{\mathrm{graph}}$ before edge
prediction.
\end{itemlist}

The reported experiments use hidden dimension 48, 4 attention heads, 2 layers,
focal binary cross-entropy ($\gamma=2.0$, $\alpha=0.25$), DropEdge regularization
(rate 0.1), a 5-epoch warmup, and cosine learning-rate decay. The edge predictor
is a small multilayer perceptron over the concatenated endpoint embeddings. The
model is intentionally compact; the gains we report come from the training
paradigm and the dual-stream design rather than from scale.

\subsection{ENCODE cCRE mapping and downstream models}

We mapped ENCODE candidate \emph{cis}-regulatory element (cCRE)
annotations\cite{encode2012,moore2020encode} to HPRC GRCh38 graph nodes by
interval overlap, yielding 303{,}425 labeled reference-path nodes across all 24
chromosomes. The implemented frozen-probe protocol uses each rotating HPRC
checkpoint only in the fold for which its test chromosomes were absent during
pretraining. Within a fold, that single checkpoint embeds downstream train,
validation, and test chromosomes so latent coordinates are not mixed across
independently trained encoders. Coordinate-only, degree-only, structural,
mono/di/tri-nucleotide-composition, frozen-embedding, sequence-plus-embedding,
training-frequency, and random baselines use the identical embedded-node
universe. Temperature and F1 threshold are fitted on validation chromosomes
only. This pipeline is implemented and tested, but its server matrix was not
complete at the evidence freeze and is not reported as an empirical result.

\subsection{Haplotype-specific methylation pilot}
\label{sec:haplotype-methods}

The exploratory methylation analysis contains 140{,}677 paired chromosome-8
H1/H2 loci from 10 HPRC donors. Nested outer splits hold out complete donors,
and every comparison uses identical paired loci. The strongest baseline uses
whole-window sequence summaries; residual models then test whether graph and
annotation features add prediction. Because the pretraining graph pools donors,
this is donor-held-out supervised evaluation but not donor-held-out
representation pretraining. The analysis is explicitly a pilot and cannot
support genome-wide haplotype-functional claims.

\section{Results}

\subsection{Verified all-chromosome execution}

The frozen server workflow completed 180 GPU jobs without recorded failure: 30
final checkpoint-production runs, 120 rotating chromosome-fold runs, 18
cross-graph transfer runs, and 12 release-transfer runs. The four parsed SV
graphs contain 389{,}065--934{,}154 segments and 562{,}222--1{,}356{,}160 links;
all fatal endpoint and orientation validation counts are zero. The rotating
folds, rather than the final all-data runs, are the primary scientific
evaluation because each canonical chromosome is held out exactly once.

\subsection{Chromosome-held-out topology reconstruction}

Table~\ref{tab:current_primary} reports means across three initialization seeds,
with intervals obtained by chromosome-level bootstrap. HPRC R2 and HGSVC3 both
show high ranking performance across all 24 chromosomes. The shared
HPRC-R2--HGSVC3 model also performs strongly when evaluated separately in each
graph, and the official integrated graph provides a consistent construction
stress test. These values establish chromosome-held-out local-topology
prediction; they do not establish donor- or haplotype-held-out learning because
each graph was parsed as pooled topology.

\begin{table}[t]
\tbl{Primary rotating chromosome-fold results. Strict induced-subgraph and
one-hop-expanded settings are separate tasks, not a matched context ablation.
Intervals resample chromosomes.}
{\small
\begin{tabular}{@{}lllccc@{}}
\toprule
Training regime & Evaluation graph & Task & AUROC & AUPRC & AUPRC 95\% CI\\
\colrule
HPRC R2 & HPRC R2 & strict & 0.9585 & 0.9501 & 0.9380--0.9529\\
HPRC R2 & HPRC R2 & one-hop & 0.9885 & 0.9880 & 0.9804--0.9891\\
HGSVC3 & HGSVC3 & strict & 0.9795 & 0.9745 & 0.9662--0.9766\\
HGSVC3 & HGSVC3 & one-hop & 0.9881 & 0.9868 & 0.9790--0.9887\\
HPRC R2 + HGSVC3 & HPRC R2 & strict & 0.9839 & 0.9797 & 0.9742--0.9822\\
HPRC R2 + HGSVC3 & HPRC R2 & one-hop & 0.9919 & 0.9916 & 0.9858--0.9923\\
HPRC R2 + HGSVC3 & HGSVC3 & strict & 0.9827 & 0.9786 & 0.9712--0.9806\\
HPRC R2 + HGSVC3 & HGSVC3 & one-hop & 0.9919 & 0.9909 & 0.9866--0.9918\\
Integrated graph & integrated graph & strict & 0.9812 & 0.9768 & 0.9708--0.9792\\
Integrated graph & integrated graph & one-hop & 0.9906 & 0.9894 & 0.9863--0.9907\\
\botrule
\end{tabular}}
\label{tab:current_primary}
\end{table}

Performance is not uniform across chromosomes. HPRC strict chrY is the clearest
stress case (AUPRC 0.8676), followed by chr18 (0.9316), chr15 (0.9336), and
chr22 (0.9350). HGSVC3 strict chrY reaches AUPRC 0.9129. These negative results
argue against extrapolating from favorable chromosomes such as chr8 and chr19.

\subsection{Context exposure, transfer, and calibration}

Expanded slices expose substantially more topology than strict slices. Across
the full-coordinate benchmarks, the expanded-to-strict ratio is 1.5869 for
visible HPRC nodes and 1.8723 for HPRC links; the corresponding HGSVC3 ratios
are 1.5731 and 1.8649. No pair among 134{,}775 interval comparisons satisfies
the preregistered joint 10\% node/link exposure caliper. Although expanded-task
scores are numerically higher, this comparison cannot distinguish better
representation learning from the easier information environment.

Cross-graph transfer is directional (Table~\ref{tab:current_transfer}). HPRC
R2-to-HGSVC3 strict transfer is weak, while its expanded-task transfer remains
high. The reverse direction is stronger for strict transfer. Four shared donors
represent 6.15\% of HGSVC3 and 1.73\% of HPRC R2, so neither direction is an
independent-donor test. Release transfer is also asymmetric and is treated as a
construction stress test.

\begin{table}[t]
\tbl{Zero-shot cross-graph and release-transfer AUPRC (mean across three seeds).}
{\begin{tabular}{@{}lcc@{}}
\toprule
Direction & Strict & One-hop\\
\colrule
HPRC R2 $\rightarrow$ HGSVC3 & 0.5997 & 0.9408\\
HGSVC3 $\rightarrow$ HPRC R2 & 0.7247 & 0.8902\\
Integrated graph $\rightarrow$ HPRC R2 & 0.7007 & 0.8836\\
HPRC R1.1 $\rightarrow$ HPRC R2 & 0.7063 & 0.7846\\
HPRC R2 $\rightarrow$ HPRC R1.1 & 0.5669 & 0.8838\\
\botrule
\end{tabular}}
\label{tab:current_transfer}
\end{table}

Raw held-out probabilities are overconfident. Validation-only temperature
scaling reduces HPRC strict ECE from 0.1799 to 0.1416 and HPRC one-hop ECE from
0.1399 to 0.0744; for HGSVC3 it reduces ECE from 0.1718 to 0.0432 (strict) and
from 0.1368 to 0.0548 (one-hop). F1 thresholds chosen exclusively on validation
chromosomes improve thresholded metrics, but ranking metrics remain primary and
the residual HPRC strict miscalibration is disclosed.

\subsection{Functional evidence remains exploratory}

No genome-wide biological downstream result is included in the frozen server
matrix. A prior chromosome-8 methylation pilot contains 140{,}677 within-donor
H1/H2 pairs from 10 donors. A strong whole-window sequence model reaches
donor-macro Spearman 0.3204. Adding graph plus annotation residuals increases
Spearman by only 0.0040 (95\% CI 0.0024--0.0055), while the graph-only residual
gain is 0.00004 (95\% CI $-0.00079$--0.00067). Graph divergence is associated
with absolute methylation difference (Spearman 0.1704), but this descriptive
association is not evidence of useful incremental prediction. A
chromosome-cross-fitted cCRE probe and phased-SV evaluation are implemented but
remain outside the executed evidence set until their server matrices complete.

\iffalse

% Author rule: all empirical performance numbers, sample counts,
% class-balance statistics, bootstrap intervals, and figure/table result
% values appear in Section 4 only. Earlier sections summarize conclusions
% qualitatively and Methods keeps protocol constants only.

\subsection{Dataset and evaluation summary}

Table~\ref{tab:results_dataset_summary} summarizes the graph and annotation
coverage used in the reported experiments. HPRC provides the primary
pretraining and internal held-out evaluation graph; HGSVC3 provides the frozen
external-transfer graph; and ENCODE cCRE intervals provide the downstream
functional labels after interval overlap with HPRC GRCh38-path nodes.

\begin{table}[t]
\tbl{Dataset and evaluation summary used in the reported experiments.}
{\small
\begin{tabular}{@{}lllrrl@{}}
\toprule
Dataset & Build/reference & Scope & Nodes & Edges & Role\\
\colrule
HPRC & v2 / GRCh38 path & canonical chrom. & 751{,}237 & 1{,}097{,}658 & training/internal test\\
HGSVC3 & CHM13 graph & chr19/21/22/Y & 72{,}033 & 100{,}546 & external transfer\\
ENCODE cCRE & GRCh38 annotation & mapped to HPRC & 303{,}425 & -- & downstream task\\
\botrule
\end{tabular}}
\label{tab:results_dataset_summary}
\end{table}

\textbf{Experimental setup.} GraphGenome-FM is pretrained with the shared,
chromosome-held-out paradigm of Section~\ref{sec:shared}. Training uses all
canonical chromosomes except chr1, chr8, chr19, and chrY (held out for testing)
and chr16 (validation); strict and one-hop closures are trained and evaluated
separately. Unless stated otherwise, link-prediction results use
distance-matched negatives. For external transfer we freeze the HPRC checkpoint
and evaluate it directly on HGSVC3 slices with no retraining or hyperparameter
tuning on external data. The downstream cCRE classifiers are trained on HPRC
nodes under the leakage-safe policy above and reported in two evaluation
universes.

\textbf{Leakage-audited evaluation.} We report unmasked message-passing runs
only as diagnostics and use query-edge-masked inference as the primary
leakage-audited estimate. This comparison directly tests how much direct edge
availability inflates link-prediction scores.

\begin{figure}[t]
\begin{center}
\includegraphics[width=0.98\textwidth]{proposed/proposed_fig1_model_architecture_user_improved.png}
\end{center}
\caption{GraphGenome-FM overview. A pangenome graph slice is represented with
node, edge, orientation, and coordinate features. A coordinate-aware attention
stream and a graph-topological attention stream are combined by a learned fusion
gate, producing edge scores for self-supervised pretraining and reusable node
embeddings for downstream tasks.}
\label{fig:study_overview}
\end{figure}

\subsection{GraphGenome-FM learns pangenome edge structure across held-out
HPRC chromosomes}

\textbf{Held-out prediction.} Trained on a subset of chromosomes and tested on
chromosomes it never saw, GraphGenome-FM predicts pangenome edge structure for
both closures (Table~\ref{tab:hprc_pretrain}). Under distance-matched negatives,
the original unmasked diagnostic is saturated (0.9846 strict, 0.9957 one-hop),
but query-edge-masked inference gives the more conservative estimate: 0.8270
strict AUROC and 0.9310 one-hop AUROC, with AUPRC 0.6959 and 0.9333,
respectively. The masked audit shows that direct-edge availability inflated the
headline AUROCs, but the model still learns signal beyond coordinate matching,
especially in one-hop graph context. \textbf{Generalization gap.} In the
unmasked training diagnostic, the near-zero gap between validation and held-out
performance indicates that the shared encoder learns recurring graph structure
rather than only fitting one chromosome.

\begin{table}[t]
\tbl{HPRC chromosome-held-out link prediction. Validation is chr16; the held-out
test chromosomes are chr1, chr8, chr19, and chrY. Query-edge-masked rows remove
positive query edges from the message-passing graph at inference time.}
{\begin{tabular}{@{}llllcc@{}}
\toprule
Benchmark & Closure & Negatives & Audit & AUROC & AUPRC\\
\colrule
HPRC base & strict & random & unmasked & 0.9889 & --\\
HPRC base & 1-hop & random & unmasked & 0.9914 & --\\
HPRC hard-negative & strict & distance-matched & unmasked & 0.9846 & --\\
HPRC hard-negative & 1-hop & distance-matched & unmasked & 0.9957 & --\\
HPRC hard-negative & strict & distance-matched & query-edge masked & 0.8270 & 0.6959\\
HPRC hard-negative & 1-hop & distance-matched & query-edge masked & 0.9310 & 0.9333\\
\botrule
\end{tabular}}
\label{tab:hprc_pretrain}
\end{table}

\subsection{Distance-matched negatives confirm the task is not driven by trivial
negatives}
\label{sec:robustness}

\textbf{Coordinate baselines collapse under distance matching.} A logistic
baseline that uses only endpoint coordinate distance reaches strict AUROC near
0.97 under random negatives but collapses to about 0.52 under distance-matched
negatives, confirming that random non-edges make the strict task artificially
easy. In the unmasked diagnostic, GraphGenome-FM instead remains strong on every
held-out chromosome (Table~\ref{tab:hprc_per_chrom}). Under query-edge masking,
the same checkpoints remain above simple topology heuristics
(Table~\ref{tab:topology_heuristics}), but the drop from 0.9846 to 0.8270 on
strict closure shows why masked inference is the primary leakage-audited
estimate.

\begin{table}[t]
\tbl{Unmasked distance-matched HPRC held-out AUROC by chromosome. Each value
averages 10 held-out windows for the corresponding closure; query-edge-masked
pooled metrics are reported in Table~\ref{tab:hprc_pretrain}.}
{\begin{tabular}{@{}lcc@{}}
\toprule
Held-out chromosome & Strict AUROC & 1-hop AUROC\\
\colrule
chr1  & 0.9887 & 0.9976\\
chr8  & 0.9819 & 0.9945\\
chr19 & 0.9769 & 0.9919\\
chrY  & 0.9911 & 0.9986\\
\botrule
\end{tabular}}
\label{tab:hprc_per_chrom}
\end{table}

\begin{table}[t]
\tbl{Topology-only heuristic baselines on the HPRC hard-negative held-out test
split. Positive query edges are removed before scoring each candidate. Values are
pooled AUROC/AUPRC over chr1, chr8, chr19, and chrY.}
{\begin{tabular}{@{}lcc@{}}
\toprule
Heuristic & Strict & 1-hop\\
\colrule
Common neighbors & 0.492/0.493 & 0.533/0.517\\
Jaccard & 0.494/0.512 & 0.533/0.523\\
Adamic--Adar & 0.491/0.488 & 0.531/0.508\\
Resource allocation & 0.491/0.488 & 0.531/0.508\\
Preferential attachment & 0.053/0.395 & 0.246/0.385\\
Shortest path & 0.440/0.473 & 0.626/0.577\\
GraphGenome-FM, masked inference & 0.827/0.696 & 0.931/0.933\\
\botrule
\end{tabular}}
\label{tab:topology_heuristics}
\end{table}

\textbf{Degree matching and seed variance.} Additional unmasked diagnostics
support the distance- and degree-control story (Fig.~\ref{fig:robustness}). When
the distance-matched negatives are also degree-matched, held-out AUROC remains
0.9839 (strict) and 0.9929 (one-hop). Across three seeds on the hard-negative
benchmark, held-out AUROC averages $0.9814\pm0.0030$ (strict) and
$0.9911\pm0.0049$ (one-hop). These unmasked controls measure robustness to
sampling choices; the query-edge-masked rows in Table~\ref{tab:hprc_pretrain}
remain the primary leakage-audited estimates.
\textbf{Non-overlapping windows.} A first non-overlapping-window run is reported
as a stress test rather than a headline result: the sampler accepted only nine
held-out slices, and performance dropped to 0.7785 (strict) and 0.9027
(one-hop). The larger strict drop suggests that overlap between sampled windows
can matter; this run is therefore useful for planning a denser tiled benchmark
but is underpowered relative to the 40-slice chromosome-held-out evaluation.

\begin{figure}[t]
\begin{center}
\includegraphics[width=0.92\textwidth]{proposed/proposed_fig11_robustness_summary.png}
\end{center}
\caption{Hard-negative and reproducibility checks for HPRC link prediction.
Distance-matched and degree-plus-distance-matched negatives, and a three-seed
rerun, all retain high held-out AUROC in the unmasked diagnostic setting. The
non-overlapping-window pilot (nine held-out slices) is an underpowered stress
test rather than a headline number, especially for strict closure.}
\label{fig:robustness}
\end{figure}

\FloatBarrier

\subsection{Ablations localize the link-prediction signal to graph topology}

A component ablation isolates the coordinate stream, the graph stream, and the
fusion gate on the distance-matched benchmark
(Table~\ref{tab:architecture_ablation}, Fig.~\ref{fig:hprc_ablation}).
These ablations were run under the original unmasked message-passing diagnostic,
so they localize signal between model components; masked ablations are listed
as future robustness work.
\textbf{Coordinate stream.} Alone it is near chance for strict edge recovery
(0.566), confirming that genomic offset features do not explain the task once
coordinate shortcuts are removed. \textbf{Graph stream.} Alone it recovers most
of the full-model performance (0.974 strict), showing that local graph topology
is the dominant signal. \textbf{Fusion gate.} Adding the coordinate stream
without a gate is strong (0.982 strict), and the learned fusion gate provides a
small but consistent additional gain (0.985 strict), which is why we retain it.
The one-hop closure is uniformly easier and saturates near 0.99 for all
topology-aware variants.

\begin{table}[t]
\tbl{Architecture ablation on the unmasked distance-matched HPRC diagnostic.
Values are held-out AUROC on chr1, chr8, chr19, and chrY, reported as
strict/1-hop.}
{\begin{tabular}{@{}llllc@{}}
\toprule
Model & Coordinate stream & Graph stream & Fusion gate & Held-out AUROC\\
\colrule
Coordinate-only & yes & no & no & 0.5658/0.8921\\
Graph-only & no & yes & no & 0.9738/0.9940\\
Dual-stream, no gate & yes & yes & no & 0.9823/0.9938\\
Full GraphGenome-FM & yes & yes & yes & 0.9846/0.9957\\
\botrule
\end{tabular}}
\label{tab:architecture_ablation}
\end{table}

\begin{figure}[t]
\begin{center}
\includegraphics[width=0.92\textwidth]{proposed/proposed_fig3_hprc_ablation_polished.png}
\end{center}
\caption{HPRC architecture ablation on held-out chromosomes in the unmasked
diagnostic setting. The coordinate stream alone is weak for strict edge recovery,
while graph-only and dual-stream models retain high performance. The full gated
GraphGenome-FM model is best overall in this diagnostic.}
\label{fig:hprc_ablation}
\end{figure}

\FloatBarrier

\subsection{The frozen HPRC checkpoint transfers to HGSVC3}

\textbf{Zero-shot transfer.} We next asked whether an HPRC-pretrained checkpoint
generalizes to an independently constructed pangenome graph. Freezing the
checkpoint and applying it to HGSVC3 without retraining, an initial chr22-only
evaluation reached strict AUROC 0.9940 and one-hop 0.9948; the expanded
evaluation with distance-matched negatives and four CHM13 targets pooled to
strict AUROC 0.9916 and one-hop 0.9955. Query-edge masking changes this
interpretation substantially (Table~\ref{tab:hgsvc_transfer}). The masked one-hop
transfer remains strong (AUROC/AUPRC 0.892/0.903), but masked strict transfer
drops near chance (0.559/0.556). Thus the HGSVC result supports portability for
one-hop neighborhood scoring, but not yet strict-edge recovery under the
leakage-audited protocol. \textbf{Per-chromosome consistency.} The masked one-hop
result is stable across all four targets, including chrY
(Table~\ref{tab:hgsvc_per_chrom}); masked strict transfer is weak for every
target. Because no HGSVC data were used for training or tuning, this is a
conservative cross-reference stress test rather than a tuned external benchmark.
\textbf{Calibration.} The reliability curves in Fig.~\ref{fig:reliability} are
based on the earlier unmasked predictions and should be interpreted as candidate
ranking diagnostics, not as final masked calibration.

\begin{table}[t]
\tbl{External HGSVC3 evaluation using frozen HPRC checkpoints. Unmasked rows are
diagnostics; query-edge-masked rows remove positive query edges before scoring
and are the leakage-audited transfer estimates.}
{\begin{tabular}{@{}llcccc@{}}
\toprule
Evaluation & Closure & Targets & AUROC & AUPRC & Brier\\
\colrule
Unmasked diagnostic & strict & chr19/21/22/Y & 0.9916 & 0.9864 & 0.0295\\
Unmasked diagnostic & 1-hop & chr19/21/22/Y & 0.9955 & 0.9940 & 0.0322\\
Query-edge masked & strict & chr19/21/22/Y & 0.5591 & 0.5555 & 0.2621\\
Query-edge masked & 1-hop & chr19/21/22/Y & 0.8924 & 0.9028 & 0.1348\\
\botrule
\end{tabular}}
\label{tab:hgsvc_transfer}
\end{table}

\begin{table}[t]
\tbl{Query-edge-masked HGSVC3 transfer by target chromosome
(distance-matched negatives). Values are mean AUROC/AUPRC across 10 slices per
target.}
{\begin{tabular}{@{}lcccc@{}}
\toprule
Target & Strict AUROC & Strict AUPRC & 1-hop AUROC & 1-hop AUPRC\\
\colrule
chr19 & 0.5949 & 0.6061 & 0.8748 & 0.8766\\
chr21 & 0.5308 & 0.5247 & 0.9059 & 0.9161\\
chr22 & 0.5434 & 0.5478 & 0.8999 & 0.9137\\
chrY  & 0.6003 & 0.5855 & 0.8762 & 0.8913\\
\botrule
\end{tabular}}
\label{tab:hgsvc_per_chrom}
\end{table}

\begin{figure}[t]
\begin{center}
\includegraphics[width=0.88\textwidth]{proposed/proposed_fig4_hgsvc_transfer_dotplot.png}
\end{center}
\caption{External transfer from HPRC to HGSVC3 under query-edge masking. Positive
query edges are removed before scoring and no HGSVC data are used for training or
tuning. One-hop AUROC/AUPRC remains high across targets, whereas strict-edge
transfer is weak under the leakage-audited protocol.}
\label{fig:hgsvc_transfer}
\end{figure}

\begin{figure}[t]
\begin{center}
\includegraphics[width=0.72\textwidth]{proposed/proposed_fig10_reliability_curves.png}
\end{center}
\caption{Reliability curves for predicted edge probabilities. HPRC held-out and
HGSVC transfer predictions track the perfect-calibration diagonal reasonably
well, with low Brier scores, indicating that the model's confidence is usable for
downstream candidate ranking.}
\label{fig:reliability}
\end{figure}

\FloatBarrier

\subsection{Application: HGSVC graph-imputation candidate scoring}

\textbf{Candidate scoring.} A portable edge-prediction model can propose missing
or weakly supported adjacencies in an existing graph. Using the frozen HPRC
one-hop checkpoint, we scored sampled HGSVC3 non-edges and ranked them as
candidate graph-enhancement edges. The unmasked HGSVC edge-prediction diagnostic
is well calibrated (AUROC 0.9955, AUPRC 0.9940, Brier 0.0322), but this
calibration should be treated as a ranking diagnostic until the candidate-scoring
pipeline is rerun and validated under the stricter masked protocol. The scoring
run produced 34{,}573 raw candidate rows that de-duplicate to 26{,}860 unique
oriented candidate pairs (Table~\ref{tab:hgsvc_imputation},
Fig.~\ref{fig:hgsvc_imputation}). \textbf{Score distribution.} Scores are
heavily concentrated near zero, with only 661 unique candidates scoring
$\geq 0.5$ and 7 scoring $\geq 0.9$, so the ranked list is a small, prioritized
set rather than a flood of low-confidence guesses. Because these candidates have
not yet been compared against a newer pangenome build, we present this as
preliminary candidate scoring; validating recovery against a newer construction
is the explicit next step.

\begin{table}[t]
\tbl{HGSVC graph-imputation candidate scoring. Candidates are sampled non-edges
scored as possible missing adjacencies. Calibration metrics are from the
unmasked edge-prediction diagnostic and should be replaced by masked/validated
metrics before making a recovery claim.}
{\begin{tabular}{@{}lr@{}}
\toprule
Metric & Value\\
\colrule
Raw candidate rows & 34,573\\
Unique oriented candidate pairs & 26,860\\
Calibration AUROC & 0.9955\\
Calibration AUPRC & 0.9940\\
Brier score & 0.0322\\
Unique candidates with score $\geq 0.5$ & 661\\
Unique candidates with score $\geq 0.8$ & 68\\
Unique candidates with score $\geq 0.9$ & 7\\
\botrule
\end{tabular}}
\label{tab:hgsvc_imputation}
\end{table}

\begin{figure}[t]
\begin{center}
\includegraphics[width=0.98\textwidth]{proposed/proposed_fig7_hgsvc_imputation_polished.png}
\end{center}
\caption{Preliminary HGSVC graph-imputation scoring. The HPRC-trained one-hop
model scores sampled HGSVC non-edges and ranks candidate missing adjacencies; the
right panel shows high-confidence candidates per target. The ranked list should
next be validated against a newer HGSVC graph build.}
\label{fig:hgsvc_imputation}
\end{figure}

\FloatBarrier

\subsection{Downstream cCRE prediction: evaluation universe and comparative
baselines}
\label{sec:ccre}

\textbf{Evaluation universe.} The cCRE experiments are leakage-controlled but
less mature than the link-prediction and transfer results, and their main lesson
is that the evaluation universe matters as much as the model class. On all
labeled held-out GRCh38 nodes, coordinate and structural logistic baselines
reach AUROC 0.837 and 0.836, and a DeepGene-style serialized Transformer reaches
AUROC 0.855 (Table~\ref{tab:comparative_baselines}); these all-node baselines
evaluate 31{,}560 nodes. The graph-native GATs currently operate on the much
smaller benchmark-window universe: as Fig.~\ref{fig:ccre_coverage} shows,
strict-or-one-hop windows cover only 1{,}330 of those held-out labeled nodes
(4.21\%). \textbf{Within-window comparison.} On this fair window universe,
coordinate, structural, and linearized logistic baselines remain below 0.60
AUROC, the serialized Transformer reaches 0.6928, and graph-native GATs reach
0.6227 from scratch, 0.6868 with a frozen pretrained encoder, and 0.6734 after
fine-tuning. We therefore make a descriptive pretraining-benefit claim within
the window universe, not a graph-native superiority claim over serialized
baselines and not an all-node superiority claim. We interpret these cCRE
differences descriptively and do not claim statistical superiority between
window models.

\begin{table}[t]
\tbl{Leakage-safe binary cCRE comparison. All cCRE models exclude
reference- and path-status features. All-node and window rows are separate
evaluation universes; the fair graph-native comparison is within the window
rows.}
{\begin{tabular}{@{}lllccc@{}}
\toprule
Method & Universe & Edge input & AUROC & AUPRC & Macro-F1\\
\colrule
Coordinate logistic & all-node & no & 0.8370 & 0.8743 & 0.7904\\
Structural logistic & all-node & feature summary & 0.8362 & 0.8746 & 0.7908\\
Linearized-graph logistic & all-node & serialized features & 0.8381 & 0.8750 & 0.7900\\
Serialized Transformer & all-node & serialized nodes & 0.8554 & 0.8894 & 0.7792\\
Structural MLP & all-node & feature summary & 0.8362 & 0.8761 & 0.7845\\
Random forest & all-node & feature summary & 0.8159 & 0.8600 & 0.7410\\
Coordinate logistic & window & no & 0.5939 & 0.4576 & 0.5591\\
Structural logistic & window & feature summary & 0.5834 & 0.4156 & 0.5681\\
Linearized-graph logistic & window & serialized features & 0.5647 & 0.3422 & 0.4844\\
Serialized Transformer & window & serialized nodes & 0.6928 & 0.5247 & 0.4003\\
GAT scratch & window & graph edges & 0.6227 & 0.4473 & 0.6419\\
GAT frozen pretrained & window & graph edges & 0.6868 & 0.5020 & 0.6518\\
GAT fine-tuned pretrained & window & graph edges & 0.6734 & 0.4937 & 0.6159\\
\botrule
\end{tabular}}
\label{tab:comparative_baselines}
\end{table}

\begin{figure}[t]
\begin{center}
\includegraphics[width=0.85\textwidth]{proposed/proposed_fig8_ccre_coverage_audit.png}
\end{center}
\caption{cCRE evaluation-universe audit. Benchmark windows cover only a small
fraction of the all-node held-out cCRE test set (strict 4.03\%, one-hop 3.60\%,
strict-or-one-hop 4.21\% of 31{,}560 nodes). This coverage gap is why all-node
baselines and window graph-native models are reported as separate universes.}
\label{fig:ccre_coverage}
\end{figure}

\begin{figure}[t]
\begin{center}
\includegraphics[width=0.98\textwidth]{proposed/proposed_fig5_ccre_fair_comparison.png}
\end{center}
\caption{Leakage-safe cCRE comparison. Binary cCRE prediction is strong for
all-node coordinate and serialized baselines and harder on the window-restricted
universe. In that universe, pretrained GATs improve over scratch and feature
baselines, while the serialized-window Transformer remains competitive and is
slightly strongest by binary AUROC in the current run. For grouped cCRE
prediction on windows, pretrained graph representations provide modest gains
over training from scratch in the descriptive comparison.}
\label{fig:ccre_binary_grouped}
\end{figure}

\textbf{Multiclass and grouped tasks.} Reduced groupings are markedly more
stable than the full label space: macro-F1 is 0.563 for three classes, 0.412 for
four, 0.337 for five, and only 0.178 for the original nine classes
(Fig.~\ref{fig:ccre_binary_grouped}). We therefore treat binary and three-class
prediction as the primary cCRE formulations and keep the four-, five-, and
nine-class results as supplementary stress tests. For the grouped GAT,
pretrained representations help most clearly in the three-class setting (0.4325
frozen versus 0.3801 from scratch) and provide smaller gains for four- and
five-class labels (Table~\ref{tab:ccre_results}). Fine-tuning is not
consistently better than freezing in the current small-window runs. These
results are a controlled downstream stress test, not evidence that graph-native
representations already solve regulatory annotation.

\begin{figure}[t]
\begin{center}
\includegraphics[width=0.98\textwidth]{proposed/proposed_fig2_ccre_pipeline_corrected.png}
\end{center}
\caption{cCRE evaluation pipeline and reporting rule. ENCODE cCRE intervals are
mapped to HPRC GRCh38-path nodes by overlap; binary, grouped, original nine-class,
and category-specific tasks are evaluated with leakage-safe feature policies.
Results are reported separately for all-node and benchmark-window-covered
universes.}
\label{fig:ccre_schematic}
\end{figure}

\begin{table}[t]
\tbl{Leakage-safe cCRE downstream summary. Binary values are AUROC; grouped and
nine-class values are macro-F1. All-node logistic and window GAT columns are
reported as separate universes.}
{\small
\begin{tabular}{@{}llcccc@{}}
\toprule
Task & Metric & All-node logit & GAT scratch & GAT frozen & GAT fine-tuned\\
\colrule
Binary & AUROC & 0.8362 & 0.6227 & 0.6868 & 0.6734\\
3-class & Macro-F1 & 0.5634 & 0.3801 & 0.4325 & 0.3941\\
4-class & Macro-F1 & 0.4122 & 0.2881 & 0.3203 & 0.2776\\
5-class & Macro-F1 & 0.3370 & 0.2201 & 0.2558 & 0.2335\\
9-class & Macro-F1 & 0.1776 & 0.1030 & 0.1357 & 0.0836\\
\botrule
\end{tabular}}
\label{tab:ccre_results}
\end{table}

\FloatBarrier

\subsection{Category-specific cCRE analysis}

Decomposing cCRE prediction into category-specific binary tasks
(Table~\ref{tab:ccre_category_breakdown},
Fig.~\ref{fig:ccre_category_performance}) is more interpretable than the
nine-class task and supports a concrete hypothesis for where graph structure
should help. \textbf{Enhancer categories.} Enhancer-associated categories are
well captured by all-node baselines (dELS AUROC 0.864, enhancer-like 0.858),
consistent with strong sequence-local signal. \textbf{Rare, topologically
situated categories.} Rarer categories are harder: TF/CTCF-associated reaches
AUROC 0.798 but AUPRC only 0.380, and promoter-like and PLS are limited
primarily by class rarity (PLS positive fraction 0.011, AUPRC 0.030). The
categories with high AUROC but low AUPRC, especially TF/CTCF-associated, are the
natural targets for a graph-topology-aware follow-up, because they combine
biological relevance with imbalance that linear baselines handle poorly.

\begin{table}[t]
\tbl{Category-specific all-node cCRE binary tasks using feature-based logistic
baselines.}
{\small
\begin{tabular}{@{}lrrrrl@{}}
\toprule
cCRE group & Positive frac. & AUROC & AUPRC & Macro-F1 & Interpretation\\
\colrule
dELS & 0.3832 & 0.8636 & 0.8612 & 0.8339 & strongest category\\
enhancer-like & 0.4209 & 0.8582 & 0.8680 & 0.8221 & strong local signal\\
pELS & 0.0953 & 0.8262 & 0.5764 & 0.7784 & good AUROC, lower AUPRC\\
TF/CTCF-assoc. & 0.0775 & 0.7984 & 0.3804 & 0.6941 & graph follow-up target\\
open chromatin & 0.1158 & 0.7840 & 0.5501 & 0.7394 & moderate\\
promoter-like & 0.0584 & 0.7460 & 0.3378 & 0.6729 & difficult but usable\\
PLS & 0.0109 & 0.6265 & 0.0303 & 0.5241 & too rare in current setup\\
\botrule
\end{tabular}}
\label{tab:ccre_category_breakdown}
\end{table}

\begin{figure}[t]
\begin{center}
\includegraphics[width=0.92\textwidth]{proposed/proposed_fig6_ccre_category_lift.png}
\end{center}
\caption{Category-specific cCRE performance against class rarity. Enhancer-like
and dELS categories are well captured by simple all-node baselines, while
promoter-like and PLS remain substantially harder because of rarity and class
imbalance. AUPRC tracks prevalence closely for the rarest categories.}
\label{fig:ccre_category_performance}
\end{figure}

\FloatBarrier

\subsection{Direct haplotype-specific methylation}
\label{sec:haplotype-results}

The exact-path cohort contains 10 HPRC donors, 125 chromosome-8 path
fragments, 284{,}375 labeled haplotype windows, and 140{,}677 reciprocal H1/H2
pairs. Every retained pair uses shared graph-node anchors; no 31-mer fallback
pair is retained. Frozen GraphGenome-FM and Nucleotide Transformer archives
contain the same 281{,}354 unique, finite embedding IDs. Exact RepeatMasker and
HMMFlagger tracks are available for nine donors; HG002 retains the graph-native
copy-number, mappability, and SV-context proxies but not those two exact tracks.

Under leave-one-donor-out evaluation, the graph-plus-context residual has the
highest rank correlation and direction accuracy (Table~\ref{tab:haplotype}).
Relative to the frozen sequence ridge, donor-macro Spearman improves by 0.0725
(paired donor-bootstrap 95\% interval 0.0617--0.0819) and direction accuracy by
0.0185 (0.0129--0.0242). The graph-only and context-only ablations lie between
the sequence baseline and the combined model in rank-correlation point
estimates, suggesting complementary signal. The same qualitative ranking
persists when complete super-populations are held out: combined Spearman is
0.0665 versus 0.0138 for sequence, and direction accuracy is 0.5214 versus
0.5048.

\begin{table}[h!]
\tbl{Haplotype-specific methylation under nested held-out evaluation. Values
are macro-averaged across held-out donors or super-populations. Context includes
exact-track differences and handcrafted graph/SV proxies.}
{\small
\begin{tabular}{@{}llrrrr@{}}
\toprule
Split & Model & Spearman & MAE & $R^2$ & Direction\\
\colrule
Donor & Sequence ridge & 0.0133 & 0.0516 & -0.0000 & 0.5038\\
Donor & Sequence + graph residual & 0.0444 & 0.0547 & -0.0718 & 0.5123\\
Donor & Sequence + context residual & 0.0552 & 0.0544 & -0.0471 & 0.5096\\
Donor & Sequence + graph/context & \textbf{0.0857} & 0.0536 & -0.0286 & \textbf{0.5223}\\
\colrule
Population & Sequence ridge & 0.0138 & 0.0517 & 0.0004 & 0.5048\\
Population & Sequence + graph/context & \textbf{0.0665} & 0.0552 & -0.0653 & \textbf{0.5214}\\
\botrule
\end{tabular}}
\label{tab:haplotype}
\end{table}

The ranking improvement does not translate into calibrated absolute
prediction. In donor folds the combined model worsens MAE by 0.0020 (95\%
interval 0.0010--0.0034 worse) and $R^2$ by 0.0286 (0.0055--0.0608 worse);
the near-zero target distribution also lets a zero predictor obtain MAE 0.0515.
The fixed gate passes donor count, Spearman gain, and direction gain, but fails
its MAE and $R^2$ requirements. We therefore do not expand to large expression,
accessibility, or 3D-contact collections in this study.

\FloatBarrier

\section{Discussion}

The central result of this work is that a compact graph-native model can learn
pangenome edge structure beyond coordinate distance, but only when evaluated
with controls that make shortcuts visible. Shared training across slices
converts many sparse local windows into a single self-supervised problem and
removes the data-starvation failure mode of per-window models
(Section~\ref{sec:shared}). Distance-matched negatives remove the coordinate
shortcut that inflates random-negative baselines. Query-edge masking then shows
how much the saturated neural diagnostic depended on direct-edge availability:
strict performance drops, while one-hop performance remains comparatively
strong. This is a more credible graph-pretraining claim than the unmasked
diagnostic alone.

The external HGSVC result is now best interpreted as a stress test rather than a
solved transfer benchmark. Without masking, HGSVC transfer appeared nearly
saturated; with positive query edges removed, one-hop transfer remains useful
across the tested external targets, whereas strict transfer is weak. This
suggests that the current checkpoint captures reusable neighborhood context, but
not enough strict-edge signal to transfer robustly across graph constructions and
reference coordinate systems. The Y chromosome remains a valuable target because
its recurrent inversions, palindromes, and segmental duplications are exactly the
kind of structure a graph-native model should eventually handle.\cite{hallast2023ychrom}

Our cCRE analysis is deliberately cautious. All downstream cCRE runs remove
reference-status features, and the comparison spans coordinate, structural,
linearized, serialized, and graph-native models. All-node coordinate and
serialized baselines are already strong for binary cCRE prediction, so cCRE
should not be presented as a simple graph-superiority benchmark. Within the
window universe, pretrained GATs improve over scratch and feature baselines, but
the serialized-window Transformer remains highly competitive. The encouraging
signals are therefore more specific: pretraining helps the grouped window tasks
modestly, reduced and category-specific tasks are far more interpretable than
the full label space, fine-tuning is not yet consistently better than freezing,
and categories with strong ranking but lower precision identify where
topology-aware modeling could plausibly add value. Because the current cCRE
input lacks nucleotide sequence tokens and quantitative assay signal, the
natural next downstream experiment is an all-node graph-native cCRE benchmark
with bootstrap intervals and richer sequence/functional features, evaluated on
the same nodes as the all-node baselines.

\subsection{Haplotype-specific functional annotation}
\label{sec:haplotype-functional}

The current cCRE experiment projects GRCh38-defined consensus annotations onto
the GRCh38 path of the pangenome graph. It therefore tests whether alternate
graph context helps predict a reference-projected label. The methylation pilot
is different: functional observations are joined directly to maternal and
paternal assembly paths, and the prediction target is a within-donor path
difference. It therefore establishes the end-to-end machinery for a genuine
haplotype-specific target rather than merely adding alternate-path context to a
GRCh38 label.

The result is encouraging for ranking but deliberately not promoted as solved
functional annotation. Graph/context residuals improve rank correlation and
allelic direction across held-out donors and populations, and the combined
model exceeds graph-only and context-only point estimates. However, absolute
error and $R^2$ worsen, indicating that the residual magnitude is poorly
calibrated. The next modeling step is validation-only residual shrinkage or a
joint rank/Huber objective under the unchanged splits and gate, together with a
longer-context sequence baseline and chromosome-level replication.

The same path interface could later predict accessibility, histone or CTCF
signal, expression, and 3D contacts from HPRC/HGSVC and 4D Nucleome
resources.\cite{encode2012,moore2020encode,dekker2017fourdn} HGSVC3
allele-specific expression remains an attractive orthogonal application because
subset-matched RNA-seq, Iso-Seq, and Hi-C are available.\cite{logsdon2025hgsvc3}
Nevertheless, the preregistered methylation gate failed its MAE and $R^2$
conditions, so we do not acquire or claim those larger modalities here. This
keeps the current positive statement narrow: graph-aware features contain
held-out ranking information, while calibrated quantitative functional
prediction remains open.

\section{Limitations and ongoing work}
\label{sec:limits}

Several limitations remain. (i) The link-prediction objective uses synthetic
negatives; distance- and degree-matched negatives are stronger controls than
random sampling, but additional local hard-negative schemes would further test
the objective. (ii) Query-edge masking has been added and used for inference
audits, but the ablation, seed-variance, degree-matched, and non-overlap neural
runs should be rerun end-to-end with query-edge masking. (iii) A first
non-overlapping-window run was underpowered because too few held-out slices were
accepted, so a denser tiled benchmark is still needed. (iv) HGSVC external
transfer currently covers a limited target set and masked strict transfer is weak;
broader external transfer requires both more targets and masked retraining. (v)
The fair cCRE graph comparison is window-aligned; an all-node graph-native cCRE
evaluation with confidence intervals remains future work. (vi) HGSVC imputation
is candidate scoring rather than validated recovery until the ranked candidates
are compared against a newer graph build. (vii) The current cCRE labels and
model inputs lack nucleotide sequence tokens and quantitative assay signals, so
confidence or signal stratification requires richer ENCODE/SCREEN metadata.
(viii) The methylation pilot covers one chromosome, reconstructs bounded graph
features from GBZ chunks, uses only 512 bp of sequence per 10-kb window, and
shows rank improvement without MAE/$R^2$ improvement. (ix) To support
reproducibility, we will release pinned software versions and training logs with
the final version.

Immediate next steps are to (a) rerun the main HPRC ablations, degree-matched
checks, seed variance, and HGSVC transfer with query-edge masking during
training and evaluation; (b) build an all-node graph-native cCRE benchmark and
report per-chromosome AUROC, AUPRC, macro-F1, balanced accuracy, and bootstrap
intervals; (c) obtain richer ENCODE/SCREEN metadata for confidence or signal
stratification; (d) compare HGSVC imputation candidates against a newer build;
(e) extend the tiled non-overlap benchmark to adequate held-out coverage; and
(f) improve methylation residual calibration under the locked donor/population
splits, add a longer-context sequence baseline, and repeat the unchanged gate
before acquiring another functional modality.

\section{Conclusion}

GraphGenome-FM provides evidence that self-supervised graph-native pretraining
can learn human pangenome graph structure beyond coordinate distance. The most
defensible current result is the leakage-audited HPRC evaluation: query-edge
masking lowers the saturated diagnostic but leaves the learned model above
topology-only heuristics, especially for one-hop graph context. HGSVC transfer is
more limited: one-hop transfer remains promising, while strict transfer requires
masked retraining and stronger validation. The same model yields a prioritized
list of candidate missing HGSVC edges for graph enhancement, but not yet
validated recovery. The cCRE experiments identify a harder downstream challenge:
converting graph pretraining into regulatory-annotation gains under leakage-safe,
biologically meaningful, same-universe evaluation, especially against strong
serialized-window baselines. This is the right status for an early graph-native
pangenome pretraining framework. The exact-path methylation cohort moves
haplotype-specific annotation from a proposal to a direct held-out-donor test:
graph/context features improve ranking and allelic direction, but not calibrated
absolute prediction. The structural benchmark is therefore mature enough to
support a cautious functional pilot, while the failed gate identifies the
specific calibration problem that must be solved before broader modality
claims.

\fi

\section{Discussion}

The executed evidence supports a focused claim: a compact encoder trained on
masked adjacency reconstruction learns recurring local pangenome topology that
generalizes across held-out chromosomes and, more variably, across graph
constructions. The full-cohort analysis changes the story in three important
ways. First, chr8 and chr19 are not representative of every chromosome; chrY is
a reproducible failure/stress chromosome. Second, strict and expanded contexts
cannot be ordered causally because exposure matching fails completely. Third,
very high within-graph ranking does not by itself demonstrate a reusable
biological foundation model.

The present scope is therefore a strong topology-pretraining conference paper,
conditional on execution of the identical-fold neural ablations and simple
baselines. A computational-genomics journal scope additionally requires the
newly implemented cross-fitted cCRE probe, phased-SV evaluation, and
donor-path-resolved graphs. The methylation data are scientifically useful as a
negative result: sequence explains substantial signal, whereas the current
topology embedding adds negligible graph-only predictive value. It should not
anchor a separate companion paper yet.

\section{Limitations and required validation}
\label{sec:limits}

The current input consists of graph topology, genomic offsets, node length,
orientation, and derived structural features; nucleotide tokens are absent.
The graphs pool donors and haplotypes before example construction, so rotating
chromosome holdout does not imply donor or haplotype holdout. HPRC R2 and HGSVC3
share four donors, and deleting those paths from a GBWT alone would not prove
that their contributions had been removed from aggregate topology. Exact
independent-donor transfer therefore requires rematerializing path-supported
graphs or rebuilding graphs without the shared assemblies.

The two context tasks differ in node and link exposure, with zero matched pairs
at the planned caliper. They must remain separately named tasks until a
target-and-exposure-matched benchmark is constructed. Probability calibration
improves after validation-only temperature scaling but remains imperfect,
especially for HPRC strict reconstruction. In the 50-kb evaluation benchmark,
256 slices are excluded because fewer than 10 default training candidates
remain (75 HPRC R2, 52 HGSVC3, 84 HPRC R1.1, and 45 integrated); complete
5-Mb generator exclusion artifacts must also accompany submission. Finally,
the full rotating-fold ablation/baseline matrix, cCRE probe, and phased-SV
evaluation are implemented but not yet executed and cannot be described as
results.

\section{Conclusion}

GraphGenome-FM currently demonstrates transferable pangenome-topology
pretraining, not a general-purpose pangenome foundation model. Its strongest
evidence is the 24-chromosome, five-fold, three-seed reconstruction analysis and
the asymmetric cross-graph/release stress tests. The evidence-safe path forward
is to complete identical-fold baselines, then require successful sequence-aware
functional and phased-SV transfer before broadening the claim or venue.

\section*{Acknowledgments}

We thank [PI Name], members of the [Lab Name] group, and our collaborators for
discussion and feedback. This work uses public resources from the Human Pangenome
Reference Consortium, HGSVC, and ENCODE. Funding and computing acknowledgments
will be finalized before submission.

\begin{thebibliography}{99}

\bibitem{liao2023hprc}
W.-W. Liao \textit{et al.}, A draft human pangenome reference,
\textit{Nature} \textbf{617}, 312--324 (2023).

\bibitem{ebert2021hgsvc}
P. Ebert \textit{et al.}, Haplotype-resolved diverse human genomes and
integrated analysis of structural variation, \textit{Science} \textbf{372},
eabf7117 (2021).

\bibitem{garrison2018vg}
E. Garrison \textit{et al.}, Variation graph toolkit improves read mapping by
representing genetic variation in the reference, \textit{Nature Biotechnology}
\textbf{36}, 875--879 (2018).

\bibitem{hallast2023ychrom}
P. Hallast, P. Ebert, M. Loftus \textit{et al.}, Assembly of 43 human Y
chromosomes reveals extensive complexity and variation, \textit{Nature}
\textbf{621}, 355--364 (2023).

\bibitem{li2020minigraph}
H. Li, X. Feng and C. Chu, The design and construction of reference pangenome
graphs with minigraph, \textit{Genome Biology} \textbf{21}, 265 (2020).

\bibitem{hickey2024minigraphcactus}
G. Hickey \textit{et al.}, Pangenome graph construction from genome alignments
with Minigraph--Cactus, \textit{Nature Biotechnology} \textbf{42}, 663--673
(2024).

\bibitem{kipf2017gcn}
T. N. Kipf and M. Welling, Semi-supervised classification with graph
convolutional networks, in \textit{International Conference on Learning
Representations} (2017).

\bibitem{velickovic2018gat}
P. Veli\v{c}kovi\'c \textit{et al.}, Graph attention networks, in
\textit{International Conference on Learning Representations} (2018).

\bibitem{ji2021dnabert}
Y. Ji, Z. Zhou, H. Liu and R. V. Davuluri, DNABERT: pre-trained Bidirectional
Encoder Representations from Transformers model for DNA-language in genome,
\textit{Bioinformatics} \textbf{37}, 2112--2120 (2021).

\bibitem{zhou2024dnabert2}
Z. Zhou, Y. Ji, W. Li, P. Dutta, R. V. Davuluri and H. Liu, DNABERT-2:
efficient foundation model and benchmark for multi-species genomes, in
\textit{International Conference on Learning Representations} (2024).

\bibitem{dallatorre2025nt}
H. Dalla-Torre \textit{et al.}, Nucleotide Transformer: building and evaluating
robust foundation models for human genomics, \textit{Nature Methods}
\textbf{22}, 287--297 (2025).

\bibitem{zhang2025deepgene}
X. Zhang, M. Yang, X. Yin, Y. Qian and F. Sun, DeepGene: an efficient foundation
model for genomics based on pan-genome graph transformer, \textit{IEEE
Transactions on Computational Biology and Bioinformatics} \textbf{22},
3143--3152 (2025).

\bibitem{mowlaei2025stici}
M. E. Mowlaei, C. Li, O. Jamialahmadi \textit{et al.}, STICI: Split-Transformer
with integrated convolutions for genotype imputation, \textit{Nature
Communications} \textbf{16}, 1218 (2025).

\bibitem{encode2012}
The ENCODE Project Consortium, An integrated encyclopedia of DNA elements in the
human genome, \textit{Nature} \textbf{489}, 57--74 (2012).

\bibitem{moore2020encode}
J. E. Moore \textit{et al.}, Expanded encyclopaedias of DNA elements in the human
and mouse genomes, \textit{Nature} \textbf{583}, 699--710 (2020).

\bibitem{logsdon2025hgsvc3}
G. A. Logsdon \textit{et al.}, Complex genetic variation in nearly complete
human genomes, \textit{Nature} \textbf{644}, 430--441 (2025).

\bibitem{dekker2017fourdn}
J. Dekker \textit{et al.}, The 4D nucleome project, \textit{Nature}
\textbf{549}, 219--226 (2017).

\end{thebibliography}

\end{document}
